% it/figures/data_pipeline/tab_pca_variance.tex — versione italiana
\begin{table}[!ht]
    \centering
    \small
    \caption[Varianza spiegata dai modi PCA verticali.]{Varianza spiegata
    da ciascuna componente principale dei profili verticali dell'AMOC. Due modi
    raggiungono il \SI{97.86}{\percent}; la coda oltre la PC2 è rumore.}
    \label{tab:pca-variance}
    \begin{tabular}{@{}lrr@{}}
        \toprule
        \textbf{Modo} & \textbf{Varianza} & \textbf{Cumulata} \\
        \midrule
        PC1 & 89.82\% & 89.82\% \\
        PC2 &  8.03\% & 97.86\% \\
        PC3 &  1.57\% & 99.42\% \\
        PC4 &  0.37\% & 99.80\% \\
        PC5 &  0.12\% & 99.91\% \\
        PC6 &  0.05\% & 99.97\% \\
        \bottomrule
    \end{tabular}
\end{table}
